%% ============================================================
%%  econcomplex --- Documentacao Tecnica
%%  Compilar com: pdflatex econcomplex_documentation.tex
%%                pdflatex econcomplex_documentation.tex  (segunda passagem)
%% ============================================================
\documentclass[12pt, a4paper]{article}

%% --- Pacotes ---------------------------------------------------
\usepackage[utf8]{inputenc}
\usepackage[T1]{fontenc}
\usepackage[brazil]{babel}
\usepackage{lmodern}
\usepackage{geometry}
\usepackage{hyperref}
\usepackage{xcolor}
\usepackage{booktabs}
\usepackage{longtable}
\usepackage{array}
\usepackage{amsmath}
\usepackage{amssymb}
\usepackage{listingsutf8}   % substitui listings com suporte completo a UTF-8
\usepackage{graphicx}
\usepackage{caption}
\usepackage{subcaption}
\usepackage{float}
\usepackage{parskip}
\usepackage{enumitem}
\usepackage{fancyhdr}
\usepackage{titlesec}
\usepackage{tcolorbox}
\usepackage{mdframed}

%% --- Geometria ------------------------------------------------
\geometry{
  a4paper,
  top=2.5cm, bottom=2.5cm,
  left=2cm, right=2cm,
  headheight=14pt
}

%% --- Cores ----------------------------------------------------
\definecolor{codebg}{RGB}{245,245,245}
\definecolor{codeframe}{RGB}{200,200,200}
\definecolor{pyblue}{RGB}{31,119,180}
\definecolor{pyorange}{RGB}{214,110,0}
\definecolor{pygray}{RGB}{100,100,100}
\definecolor{tblhead}{RGB}{40,70,120}
\definecolor{noteblue}{RGB}{220,235,250}
\definecolor{noteborder}{RGB}{70,130,180}
\definecolor{warnbg}{RGB}{255,248,220}
\definecolor{warnborder}{RGB}{200,150,0}

%% --- Mapeamento UTF-8 para lstlistings ------------------------
%% (necessario para acentos dentro dos blocos de codigo)
\lstset{
  language=Python,
  basicstyle=\ttfamily\footnotesize,
  backgroundcolor=\color{codebg},
  frame=single,
  framerule=0.5pt,
  rulecolor=\color{codeframe},
  keywordstyle=\bfseries\color{pyblue},
  commentstyle=\itshape\color{pygray},
  stringstyle=\color{pyorange},
  numbers=left,
  numberstyle=\tiny\color{pygray},
  numbersep=8pt,
  showstringspaces=false,
  breaklines=true,
  breakatwhitespace=true,
  tabsize=4,
  captionpos=b,
  xleftmargin=1.5em,
  xrightmargin=0.5em,
  inputencoding=utf8,
  extendedchars=true,
  literate=
    {\'a}{{\'{a}}}1 {\'e}{{\'{e}}}1 {\'i}{{\'{i}}}1 {\'o}{{\'{o}}}1 {\'u}{{\'{u}}}1
    {\^a}{{\^{a}}}1 {\^e}{{\^{e}}}1 {\^o}{{\^{o}}}1
    {\~a}{{\~{a}}}1 {\~o}{{\~{o}}}1
    {\c c}{{\c{c}}}1
    {a\'}{{\'{a}}}1 {e\'}{{\'{e}}}1
    {i\'}{{\'{i}}}1 {o\'}{{\'{o}}}1 {u\'}{{\'{u}}}1
    {a\^}{{\^{a}}}1 {e\^}{{\^{e}}}1 {o\^}{{\^{o}}}1
    {a\~}{{\~{a}}}1 {o\~}{{\~{o}}}1
    {,}{{,}}1 {;}{{;}}1 {:}{{:}}1,
}

%% --- Estilo Bash (sem auto-referencia) -----------------------
\lstdefinestyle{bashstyle}{
  language=bash,
  basicstyle=\ttfamily\footnotesize,
  backgroundcolor=\color{codebg},
  frame=single,
  framerule=0.5pt,
  rulecolor=\color{codeframe},
  keywordstyle=\bfseries\color{pygray},
  commentstyle=\itshape\color{pygray},
  stringstyle=\color{pyorange},
  numbers=left,
  numberstyle=\tiny\color{pygray},
  numbersep=8pt,
  showstringspaces=false,
  breaklines=true,
  tabsize=2,
  captionpos=b,
  xleftmargin=1.5em,
  xrightmargin=0.5em,
  inputencoding=utf8,
  extendedchars=true,
  morekeywords={pip,pip3,python3,cd,jupyter,echo,source,export,conda,activate},
}

%% --- Caixas de destaque --------------------------------------
\tcbuselibrary{skins, breakable}

\newtcolorbox{notebox}{
  colback=noteblue, colframe=noteborder,
  boxrule=0.8pt, arc=3pt,
  left=6pt, right=6pt, top=4pt, bottom=4pt,
  fonttitle=\bfseries, title=Nota,
}

\newtcolorbox{warnbox}{
  colback=warnbg, colframe=warnborder,
  boxrule=0.8pt, arc=3pt,
  left=6pt, right=6pt, top=4pt, bottom=4pt,
  fonttitle=\bfseries, title={Aten\c{c}\~{a}o},
}

%% --- Cabecalho / Rodape --------------------------------------
\pagestyle{fancy}
\fancyhf{}
\fancyhead[L]{\small\texttt{econcomplex}}
\fancyhead[R]{\small Documenta\c{c}\~{a}o T\'{e}cnica}
\fancyfoot[C]{\small\thepage}
\renewcommand{\headrulewidth}{0.4pt}

%% --- Secoes ---------------------------------------------------
\titleformat{\section}{\large\bfseries\color{tblhead}}{}{0em}{}[\titlerule]
\titleformat{\subsection}{\normalsize\bfseries\color{pyblue}}{}{0em}{}
\titleformat{\subsubsection}{\normalsize\itshape}{}{0em}{}

%% --- Hyperref -------------------------------------------------
\hypersetup{
  colorlinks=true,
  linkcolor=pyblue,
  urlcolor=pyblue,
  citecolor=pygray,
  pdftitle={econcomplex -- Documentacao Tecnica},
  pdfauthor={econcomplex contributors},
}

%% --- Macros ---------------------------------------------------
\newcommand{\pkg}[1]{\texttt{\textbf{#1}}}
\newcommand{\fn}[1]{\texttt{#1}}
\newcommand{\argm}[1]{\textit{\texttt{#1}}}
\newcommand{\mat}[1]{\mathbf{#1}}

%% ============================================================
\begin{document}

%% --- Capa ----------------------------------------------------
\begin{titlepage}
  \centering
  \vspace*{3cm}
  {\Huge\bfseries\color{tblhead} econcomplex}\par
  \vspace{0.5cm}
  {\large\bfseries Documenta\c{c}\~{a}o T\'{e}cnica --- Vers\~{a}o 0.1.0}\par
  \vspace{1cm}
  \rule{10cm}{0.8pt}\par
  \vspace{0.8cm}
  {\large Biblioteca Python para Indicadores de\\
  Complexidade Econ\^{o}mica e Ci\^{e}ncia Regional}\par
  \vspace{2cm}
  {\normalsize
    \textit{Elton Freitas}
  }\par
  \vspace{3cm}
  {\small\color{pygray}
    Depend\^{e}ncias: \texttt{numpy $\geq$ 1.21}, \texttt{pandas $\geq$ 1.3}, \texttt{scipy $\geq$ 1.7}\\[4pt]
    Compat\'{i}vel com Python 3.9+
  }\par
  \vfill
  {\small\color{pygray} Mar\c{c}o de 2026}
\end{titlepage}

%% --- Indice --------------------------------------------------
\tableofcontents
\newpage

%% ============================================================
\section{Vis\~{a}o Geral}

\pkg{econcomplex} \'{e} uma biblioteca Python que consolida as melhores
implementa\c{c}\~{o}es de indicadores de Complexidade Econ\^{o}mica e de
Ci\^{e}ncia Regional, reunindo funcionalidades de quatro pacotes de
refer\^{e}ncia na literatura:

\begin{itemize}[leftmargin=1.5em]
  \item \textbf{EconGeo} (Balland, 2017) --- pacote R com foco em indicadores
        geogr\'{a}ficos e de especializa\c{c}\~{a}o regional
  \item \textbf{economiccomplexity} (Vargas Sep\'{u}lveda, 2020) --- pacote R
        com backend C++ (Armadillo) para indicadores de complexidade
  \item \textbf{py-ecomplexity} (The Growth Lab at Harvard, 2020) --- pipeline
        completo em Python com suporte a s\'{e}ries temporais
  \item \textbf{py-economic-complexity} --- implementa\c{c}\~{a}o modular em
        Python com suporte a Polars e indicadores de cross-space
\end{itemize}

\subsection{Instala\c{c}\~{a}o}

\begin{lstlisting}[style=bashstyle,
  caption={Instala\c{c}\~{a}o via pip (desenvolvimento local)}]
# A partir do diretorio raiz da biblioteca
pip install -e .

# Ou adicionando o caminho manualmente
import sys
sys.path.insert(0, "/caminho/para/econcomplex")
import econcomplex as ec
\end{lstlisting}

\subsection{Estrutura de M\'{o}dulos}

\begin{center}
\begin{tabular}{lp{9cm}}
\toprule
\textbf{M\'{o}dulo} & \textbf{Conte\'{u}do} \\
\midrule
\texttt{core}           & RCA, RPOP, Mcp, diversidade, ubiquidade, utilit\'{a}rios \\
\texttt{complexity}     & ECI/PCI (3 m\'{e}todos), ECI subnacional \\
\texttt{relatedness}    & Proximidade, densidade, dist\^{a}ncia, co-ocorr\^{e}ncia, cross-space \\
\texttt{specialization} & LQ, Hachman, Krugman, coef.\ especializa\c{c}\~{a}o, similaridade \\
\texttt{inequality}     & Gini, Gini locacional, Hoover-Gini, HHI, entropia Shannon \\
\texttt{productivity}   & PRODY, EXPY, PGI, PEII \\
\texttt{patents}        & Ease of Recombination, Modular Complexity \\
\texttt{dynamics}       & Taxas de crescimento, rastreamento de entrada/sa\'{i}da \\
\texttt{outlook}        & COI, COG (Complexity Outlook) \\
\texttt{pipeline}       & Fun\c{c}\~{a}o \texttt{compute\_complexity()} --- pipeline unificado \\
\bottomrule
\end{tabular}
\end{center}

%% ============================================================
\section{Organiza\c{c}\~{a}o dos Dados}

Esta se\c{c}\~{a}o descreve o formato de entrada esperado pela biblioteca,
garantindo que os dados estejam corretamente estruturados antes da
execu\c{c}\~{a}o de qualquer indicador.

% -- -- -- -- -- -- -- -- -- -- -- -- -- -- -- -- -- -- -- --
\subsection{Formato Padr\~{a}o: Tabela Longa (\textit{long format})}

A biblioteca \pkg{econcomplex} trabalha exclusivamente com
\textbf{tabelas no formato longo} (\textit{tidy data}),
onde cada linha representa \textbf{uma combina\c{c}\~{a}o \'{u}nica} de
regi\~{a}o $\times$ atividade (e per\'{i}odo, quando houver):

\begin{center}
\renewcommand{\arraystretch}{1.3}
\begin{tabular}{llll}
\toprule
\textbf{regiao} & \textbf{atividade} & \textbf{valor} & \textbf{ano (opcional)} \\
\midrule
\texttt{SP}  & \texttt{cnae\_10} & 12\,345 & 2022 \\
\texttt{SP}  & \texttt{cnae\_25} &  6\,789 & 2022 \\
\texttt{RJ}  & \texttt{cnae\_10} &  9\,012 & 2022 \\
\texttt{RJ}  & \texttt{cnae\_25} &  3\,456 & 2022 \\
\texttt{SP}  & \texttt{cnae\_10} & 13\,200 & 2023 \\
$\vdots$ & $\vdots$ & $\vdots$ & $\vdots$ \\
\bottomrule
\end{tabular}
\end{center}

\begin{warnbox}
  A biblioteca \textbf{n\~{a}o} aceita matrizes pivotadas (regi\~{o}es nas
  linhas, atividades nas colunas). Use \fn{pivot\_to\_matrix()} apenas para
  fun\c{c}\~{o}es de baixo n\'{i}vel que exijam a matriz explicitamente.
\end{warnbox}

% -- -- -- -- -- -- -- -- -- -- -- -- -- -- -- -- -- -- -- --
\subsection{Regras de Qualidade dos Dados}

\begin{enumerate}[leftmargin=1.5em]
  \item \textbf{Sem linhas duplicadas:} cada par (\textit{regiao},
        \textit{atividade}) deve aparecer exatamente uma vez por per\'{i}odo.
  \item \textbf{Valores n\~{a}o-negativos:} a coluna de valor deve conter
        apenas n\'{u}meros $\geq 0$.
  \item \textbf{Sem \texttt{NaN}:} remova ou impute valores ausentes antes
        de passar o DataFrame.
  \item \textbf{Granularidade consistente:} todas as regi\~{o}es devem usar
        o mesmo n\'{i}vel geogr\'{a}fico (ex.: todos munic\'{i}pios
        \textit{ou} todos estados).
  \item \textbf{Atividades homog\^{e}neas:} use uma \'{u}nica classifica\c{c}\~{a}o
        por an\'{a}lise (ex.: CNAE 2 d\'{i}gitos \textit{ou} CNAE 4
        d\'{i}gitos, n\~{a}o misturados).
\end{enumerate}

% -- -- -- -- -- -- -- -- -- -- -- -- -- -- -- -- -- -- -- --
\subsection{Nomes das Colunas --- Dicion\'{a}rio \texttt{cols}}

O pipeline usa um dicion\'{a}rio \argm{cols} para mapear os nomes reais
das colunas do seu DataFrame:

\begin{center}
\renewcommand{\arraystretch}{1.3}
\begin{tabular}{lll}
\toprule
\textbf{Chave} & \textbf{Tipo} & \textbf{Descri\c{c}\~{a}o} \\
\midrule
\texttt{"loc"}  & obrigat\'{o}ria & Coluna de regi\~{a}o / localiza\c{c}\~{a}o \\
\texttt{"act"}  & obrigat\'{o}ria & Coluna de atividade / produto / tecnologia \\
\texttt{"val"}  & obrigat\'{o}ria & Coluna de valor num\'{e}rico (empregos, exporta\c{c}\~{o}es \ldots) \\
\texttt{"time"} & opcional        & Coluna de per\'{i}odo (ano, trimestre \ldots) \\
\bottomrule
\end{tabular}
\end{center}

\begin{lstlisting}[caption={Verifica\c{c}\~{a}o e prepara\c{c}\~{a}o do DataFrame antes do pipeline}]
import pandas as pd
import econcomplex as ec

df = pd.read_csv("meus_dados.csv")

# Verificar estrutura
print(df.dtypes)
print(df.shape)
print(df.isnull().sum())        # checar NaN
print(df.duplicated().sum())    # checar duplicatas

# Limpeza basica
df = df.dropna(subset=["regiao", "cnae", "empregos"])
df = df[df["empregos"] >= 0]    # remover valores negativos
df = df.drop_duplicates(subset=["regiao", "cnae", "ano"])

# Rodar o pipeline
resultado = ec.compute_complexity(
    df,
    cols={"loc": "regiao", "act": "cnae",
          "val": "empregos", "time": "ano"},
)
\end{lstlisting}

% -- -- -- -- -- -- -- -- -- -- -- -- -- -- -- -- -- -- -- --
\subsection{Formatos de Arquivo Suportados}

\begin{center}
\renewcommand{\arraystretch}{1.3}
\begin{tabular}{lll}
\toprule
\textbf{Formato} & \textbf{Leitura pandas} & \textbf{Observa\c{c}\~{a}o} \\
\midrule
\texttt{.csv}     & \texttt{pd.read\_csv()}     & Mais comum; use \texttt{sep=";"} para CSVs BR \\
\texttt{.parquet} & \texttt{pd.read\_parquet()} & Ideal para bases grandes (RAIS, COMEX) \\
\texttt{.xlsx}    & \texttt{pd.read\_excel()}   & Requer \texttt{openpyxl}; evitar $>$100\,000 linhas \\
\texttt{.feather} & \texttt{pd.read\_feather()} & Alta performance; requer \texttt{pyarrow} \\
\texttt{.dta}     & \texttt{pd.read\_stata()}   & Arquivos Stata; preserva labels \\
\bottomrule
\end{tabular}
\end{center}

\begin{lstlisting}[caption={Leitura de diferentes formatos de arquivo}]
import pandas as pd

# CSV com separador ponto-e-virgula (padrao brasileiro)
df = pd.read_csv("rais.csv", sep=";", encoding="latin-1",
                 usecols=["municipio", "cnae2d", "ano", "vinculos"],
                 dtype={"municipio": str, "cnae2d": str})

# Parquet (preferencial para RAIS/COMEX)
df = pd.read_parquet("comex_ncm.parquet",
                     columns=["co_mun", "co_ncm", "ano", "vl_fob"])

# Excel
df = pd.read_excel("dados.xlsx", sheet_name="Sheet1",
                   usecols=["uf", "setor", "ano", "empregos"])
\end{lstlisting}

% -- -- -- -- -- -- -- -- -- -- -- -- -- -- -- -- -- -- -- --
\subsection{Exemplos de Bases Compat\'{i}veis}

\begin{description}[leftmargin=1.5em, style=nextline]
  \item[\textbf{RAIS/MTE}]
    Empregos formais por munic\'{i}pio $\times$ CNAE (2 ou 4 d\'{i}gitos)
    $\times$ ano. Coluna de valor: v\'{i}nculos ativos em 31/12.
  \item[\textbf{COMEX/MDIC}]
    Exporta\c{c}\~{o}es por munic\'{i}pio $\times$ NCM $\times$ ano.
    Coluna de valor: \texttt{VL\_FOB} (valor em US\$).
  \item[\textbf{INPI --- Patentes}]
    Pedidos de patente por Unidade da Federa\c{c}\~{a}o $\times$
    classifica\c{c}\~{a}o IPC $\times$ ano.
  \item[\textbf{Atlas da Complexidade (Harvard)}]
    Exporta\c{c}\~{o}es por pa\'{i}s $\times$ produto (HS) $\times$ ano.
    Dispon\'{i}vel em \url{https://atlas.cid.harvard.edu/}.
  \item[\textbf{Dados simulados (testes)}]
    Use \fn{ec.make\_sample\_data(n\_locs=20, n\_acts=50, seed=42)}
    para gerar um DataFrame sint\'{e}tico.
\end{description}

%% ============================================================
\section{Como Executar os Scripts}

Esta se\c{c}\~{a}o descreve todas as formas de executar a biblioteca
\pkg{econcomplex}, desde a execu\c{c}\~{a}o direta de scripts Python
at\'{e} o uso em ambientes interativos como Jupyter e IDEs.

% -- -- -- -- -- -- -- -- -- -- -- -- -- -- -- -- -- -- -- --
\subsection{Pr\'{e}-requisitos}

\begin{warnbox}
  Certifique-se de que as depend\^{e}ncias est\~{a}o instaladas antes de
  executar qualquer script.
\end{warnbox}

\begin{lstlisting}[style=bashstyle,
  caption={Instala\c{c}\~{a}o das depend\^{e}ncias}]
# Instalar dependencias principais
pip install numpy pandas scipy

# Instalar a biblioteca em modo de desenvolvimento
pip install -e /caminho/para/econcomplex

# Verificar instalacao
python3 -c "import econcomplex as ec; print(ec.__version__)"
\end{lstlisting}

% -- -- -- -- -- -- -- -- -- -- -- -- -- -- -- -- -- -- -- --
\subsection{Execu\c{c}\~{a}o Direta via Terminal}

A forma mais simples de usar a biblioteca \'{e} executar um script
\texttt{.py} diretamente no terminal:

\begin{lstlisting}[style=bashstyle,
  caption={Execu\c{c}\~{a}o de scripts no terminal}]
# Navegue ate a pasta do projeto
cd "/caminho/para/1 - Pos-doc research/data/script"

# Executar o script de exemplo completo
python3 econcomplex/exemplo_uso.py

# Executar um script proprio
python3 meu_script.py
\end{lstlisting}

Caso a biblioteca \textbf{n\~{a}o} esteja instalada via \texttt{pip},
adicione o caminho manualmente no in\'{i}cio de cada script:

\begin{lstlisting}[caption={Adicionar caminho ao sys.path (sem pip install)}]
import sys
sys.path.insert(0, "/caminho/para/econcomplex")

import econcomplex as ec
print(ec.__version__)   # 0.1.0
\end{lstlisting}

% -- -- -- -- -- -- -- -- -- -- -- -- -- -- -- -- -- -- -- --
\subsection{Execu\c{c}\~{a}o em Jupyter Notebook / JupyterLab}

\begin{lstlisting}[style=bashstyle,
  caption={Abrir Jupyter a partir da pasta de scripts}]
cd "/caminho/para/1 - Pos-doc research/data/script"
jupyter notebook
# ou
jupyter lab
\end{lstlisting}

Dentro do notebook, a primeira c\'{e}lula deve configurar o caminho:

\begin{lstlisting}[caption={C\'{e}lula de configura\c{c}\~{a}o no Jupyter}]
# Celula 1 -- configuracao
import sys, os

BASE = os.path.expanduser(
    "~/Library/CloudStorage/"
    "OneDrive-Personnel/1 - Pos-doc research/data/script/econcomplex"
)
sys.path.insert(0, BASE)

import econcomplex as ec
import numpy as np
import pandas as pd

print("econcomplex", ec.__version__, "carregado com sucesso.")
\end{lstlisting}

\begin{lstlisting}[caption={C\'{e}lulas subsequentes --- uso normal}]
# Celula 2 -- carregar dados
df = pd.read_csv("meus_dados.csv")   # colunas: regiao, setor, empregos

# Celula 3 -- calcular complexidade
resultado = ec.compute_complexity(
    df,
    cols={"loc": "regiao", "act": "setor", "val": "empregos"},
    method="eigenvector",
    compute_coi_cog=True,
)
resultado.head()

# Celula 4 -- visualizar ECI por regiao
eci = (resultado[["regiao", "eci"]]
       .drop_duplicates()
       .set_index("regiao")
       .sort_values("eci", ascending=False))
print(eci)
\end{lstlisting}

% -- -- -- -- -- -- -- -- -- -- -- -- -- -- -- -- -- -- -- --
\subsection{Execu\c{c}\~{a}o no VS Code / PyCharm / Spyder}

Em qualquer IDE, configure o \textbf{interpretador Python} e o
\textbf{diret\'{o}rio de trabalho} apontando para a pasta de scripts.

\begin{lstlisting}[style=bashstyle,
  caption={Configurar caminho no VS Code (settings.json)}]
// .vscode/settings.json
{
  "python.defaultInterpreterPath":
      "/Library/Frameworks/Python.framework/Versions/3.12/bin/python3",
  "terminal.integrated.cwd": "${workspaceFolder}"
}
\end{lstlisting}

No \textbf{Spyder}, v\'{a} em \textit{Tools $\to$ PYTHONPATH Manager}
e adicione o diret\'{o}rio \texttt{econcomplex/}.

% -- -- -- -- -- -- -- -- -- -- -- -- -- -- -- -- -- -- -- --
\subsection{Script M\'{i}nimo de Execu\c{c}\~{a}o}

Template pronto para uso com qualquer base em formato CSV:

\begin{lstlisting}[caption={Template m\'{i}nimo --- copiar e adaptar}]
"""
Template de uso da biblioteca econcomplex.
Adapte os nomes das colunas e o caminho do arquivo CSV.
"""
import sys
sys.path.insert(0, "/caminho/para/econcomplex")

import pandas as pd
import econcomplex as ec

# 1. Carregar dados
df = pd.read_csv("dados.csv")
print("Colunas:", df.columns.tolist())
print("Shape:", df.shape)

# 2. Calcular todos os indicadores
resultado = ec.compute_complexity(
    df,
    cols={
        "loc":  "nome_coluna_regiao",    # ex: "municipio", "uf"
        "act":  "nome_coluna_atividade", # ex: "cnae", "produto"
        "val":  "nome_coluna_valor",     # ex: "empregos", "exportacoes"
        "time": "nome_coluna_ano",       # opcional
    },
    method="eigenvector",   # 'eigenvector' | 'reflections' | 'fitness'
    threshold=1.0,
    compute_coi_cog=True,
)

# 3. Exportar resultados
resultado.to_csv("resultados_complexidade.csv", index=False)
print("Resultado salvo em resultados_complexidade.csv")
print("Colunas geradas:", sorted(resultado.columns.tolist()))

# 4. Indicadores individuais (opcional)
mat = ec.pivot_to_matrix(df,
                         "nome_coluna_regiao",
                         "nome_coluna_atividade",
                         "nome_coluna_valor")

eci, pci  = ec.eci_pci(mat)
hachman   = ec.hachman_index(mat)
krugman   = ec.krugman_index(mat)
entropia  = ec.shannon_entropy(mat)

saida = pd.DataFrame({
    "eci": eci, "hachman": hachman,
    "krugman": krugman, "entropia": entropia,
})
print(saida.round(3))
\end{lstlisting}

% -- -- -- -- -- -- -- -- -- -- -- -- -- -- -- -- -- -- -- --
\subsection{Execu\c{c}\~{a}o com Dados da RAIS / CNAE (caso de uso real)}

\begin{lstlisting}[caption={Uso com dados RAIS --- munic\'{i}pios x CNAE}]
import sys
sys.path.insert(0, "/caminho/para/econcomplex")

import pandas as pd
import econcomplex as ec

# Carregar RAIS agregada
# Esperado: municipio, cnae2d, ano, vinculos
rais = pd.read_parquet("rais_municipio_cnae.parquet")

# Filtrar um unico ano para analise
rais_2022 = rais[rais["ano"] == 2022].copy()

# Pipeline
resultado = ec.compute_complexity(
    rais_2022,
    cols={"loc": "municipio", "act": "cnae2d", "val": "vinculos"},
    method="eigenvector",
    compute_coi_cog=False,   # desligar COG para bases grandes
)

# Top 20 municipios mais complexos
top20 = (resultado[["municipio", "eci"]]
         .drop_duplicates()
         .sort_values("eci", ascending=False)
         .head(20))
print(top20.to_string(index=False))

resultado.to_csv("complexidade_municipios_2022.csv", index=False)
\end{lstlisting}

% -- -- -- -- -- -- -- -- -- -- -- -- -- -- -- -- -- -- -- --
\subsection{Execu\c{c}\~{a}o em Painel Multitemporal}

Para analisar a evolu\c{c}\~{a}o dos indicadores ao longo do tempo,
basta incluir a coluna temporal no dicion\'{a}rio \argm{cols}:

\begin{lstlisting}[caption={Execu\c{c}\~{a}o em painel --- m\'{u}ltiplos anos}]
# rais com coluna "ano" contendo 2010, 2015, 2020
resultado_painel = ec.compute_complexity(
    rais,
    cols={
        "loc":  "municipio",
        "act":  "cnae2d",
        "val":  "vinculos",
        "time": "ano",       # ativa o loop temporal
    },
    method="reflections",
    compute_coi_cog=False,
)

# O resultado ja contem todos os anos
print(resultado_painel["ano"].value_counts().sort_index())

# Evolucao do ECI mediano por ano
eci_evolucao = (resultado_painel[["ano", "municipio", "eci"]]
                .drop_duplicates()
                .groupby("ano")["eci"]
                .median()
                .round(4))
print(eci_evolucao)
\end{lstlisting}

\begin{notebox}
  O pipeline itera automaticamente sobre cada per\'{i}odo, garantindo que
  a matriz de proximidade e os vetores de complexidade sejam recalculados
  de forma independente para cada ano.
\end{notebox}

%% ============================================================
\section{Indicadores do N\'{u}cleo (\texttt{core})}

\subsection{Vantagem Comparativa Revelada (RCA)}

A RCA de Balassa \'{e} o principal filtro de especializa\c{c}\~{a}o:

\[
  \text{RCA}_{il} = \frac{x_{il} / \sum_l x_{il}}
                         {\sum_i x_{il} / \sum_{il} x_{il}}
\]

\begin{lstlisting}[caption={RCA cont\'{i}nuo e bin\'{a}rio}]
import econcomplex as ec
import pandas as pd

df = pd.read_csv("dados.csv")
mat = ec.pivot_to_matrix(df, "regiao", "setor", "empregos")

rca      = ec.rca(mat)               # valores continuos
mcp      = ec.mcp(mat, threshold=1)  # 1 se RCA >= 1, 0 caso contrario
rca_rpop = ec.rpop(mat)              # RCA de Puga & Turrini (continuo)
\end{lstlisting}

\subsection{Diversidade e Ubiquidade}

\begin{lstlisting}[caption={Diversidade e ubiquidade}]
div = ec.diversity(mat)    # numero de setores com RCA >= 1 por regiao
ubi = ec.ubiquity(mat)     # numero de regioes com RCA >= 1 por setor
\end{lstlisting}

%% ============================================================
\section{Complexidade Econ\^{o}mica (\texttt{complexity})}

\subsection{M\'{e}todo do Vetor Pr\'{o}prio (Hidalgo \& Hausmann, 2009)}

\[
  \text{ECI} = \frac{\vec{v}_2 - \mu(\vec{v}_2)}{\sigma(\vec{v}_2)}, \quad
  \text{PCI} = \frac{\vec{u}_2 - \mu(\vec{u}_2)}{\sigma(\vec{u}_2)}
\]
onde $\vec{v}_2$ \'{e} o segundo autovetor esquerdo da matriz $\tilde{M}$:
\[
  \tilde{M}_{ij} = \frac{M_{ij}}{k_{c,0} \cdot k_{p,0}}
\]

\begin{lstlisting}[caption={ECI e PCI pelo m\'{e}todo do vetor pr\'{o}prio}]
eci, pci = ec.eci_pci(mat, method="eigenvector")
print(eci.head())   # Series indexada por regiao
print(pci.head())   # Series indexada por produto/setor
\end{lstlisting}

\subsection{M\'{e}todo de Reflex\~{o}es (Hidalgo \& Hausmann, 2009)}

\begin{lstlisting}[caption={M\'{e}todo de Reflex\~{o}es}]
eci_r, pci_r = ec.eci_pci(mat, method="reflections", iterations=20)
\end{lstlisting}

\subsection{Fitness-Complexity (Tacchini et al., 2012)}

\[
  F_c^{(n+1)} = \sum_p M_{cp} Q_p^{(n)}, \quad
  Q_p^{(n+1)} = \frac{1}{\sum_c M_{cp} / F_c^{(n)}}
\]

\begin{lstlisting}[caption={Fitness-Complexity}]
fit, comp = ec.eci_pci(mat, method="fitness", iterations=100)
\end{lstlisting}

\subsection{ECI Subnacional}

\begin{lstlisting}[caption={ECI subnacional com PCI externo}]
# pci_nacional: PCI calculado com dados nacionais/internacionais
eci_sub = ec.subnational_eci(mat_sub, pci_nacional)
\end{lstlisting}

%% ============================================================
\section{Product Space e Relatedness (\texttt{relatedness})}

\subsection{Proximidade ($\phi$)}

\[
  \phi_{ij} = \min\!\left(
    P(\text{RCA}_i \geq 1 \mid \text{RCA}_j \geq 1),\;
    P(\text{RCA}_j \geq 1 \mid \text{RCA}_i \geq 1)
  \right)
\]

\begin{lstlisting}[caption={Proximidade discreta e cont\'{i}nua}]
phi_d = ec.proximity(mat, continuous=False)   # baseada em Mcp binario
phi_c = ec.proximity(mat, continuous=True)    # baseada em RCA continuo
\end{lstlisting}

\subsection{Densidade e Dist\^{a}ncia}

\[
  \omega_{ij} = \frac{\sum_k M_{kj} \phi_{ij}}{\sum_k \phi_{ij}}, \quad
  d_{ij} = 1 - \omega_{ij}
\]

\begin{lstlisting}[caption={Densidade e dist\^{a}ncia}]
density  = ec.density(mat, phi_d)   # DataFrame loc x act
distance = ec.distance(mat, phi_d)  # DataFrame loc x act
\end{lstlisting}

\subsection{Co-ocorr\^{e}ncia e \'{I}ndices de Relatedness}

\begin{lstlisting}[caption={Co-ocorr\^{e}ncia e \'{i}ndices de relatedness}]
cooc    = ec.co_occurrence(mat)
phi_cos = ec.cosine_proximity(mat)   # baseada em cosseno
phi_cor = ec.correlation_proximity(mat)
rca_rel = ec.relatedness(mat, phi_d)  # densidades de relatedness
\end{lstlisting}

\subsection{Cross-Space (relatedness entre classifica\c{c}\~{o}es distintas)}

\begin{lstlisting}[caption={Proximidade e relatedness cross-space}]
# phi_cross: proximidade entre setores e tecnologias (matrizes diferentes)
phi_cross = ec.cross_space_proximity(mat_setor, mat_tech)
density_c = ec.density(mat_setor, phi_cross)
\end{lstlisting}

%% ============================================================
\section{Especializa\c{c}\~{a}o Regional (\texttt{specialization})}

\begin{lstlisting}[caption={\'{I}ndices de especializa\c{c}\~{a}o}]
lq        = ec.location_quotient(mat)    # equivalente ao RCA
hachman   = ec.hachman_index(mat)        # 0 (diversificado) a 1 (especializado)
krugman   = ec.krugman_index(mat)        # diferenca estrutural entre regioes
coef_esp  = ec.spec_coefficient(mat)     # coeficiente de especializacao
sim_exp   = ec.export_similarity(mat)    # similaridade de estrutura produtiva
\end{lstlisting}

%% ============================================================
\section{Desigualdade e Concentra\c{c}\~{a}o (\texttt{inequality})}

\begin{lstlisting}[caption={Indicadores de desigualdade}]
gini_r   = ec.gini(mat)                # Gini da distribuicao do valor
gini_loc = ec.locational_gini(mat)     # Gini locacional (por setor)
hoover   = ec.hoover_gini(mat)         # Indice de Hoover
hhi      = ec.hhi(mat)                 # Herfindahl-Hirschman
entropia = ec.shannon_entropy(mat)     # entropia de Shannon (por regiao)
\end{lstlisting}

%% ============================================================
\section{Produtividade (\texttt{productivity})}

\begin{lstlisting}[caption={Produtividade e desigualdade dos produtos}]
prody = ec.prody(mat, gdp_per_capita)  # renda associada ao produto
expy  = ec.expy(mat, prody)            # sofisticacao da cesta exportadora
pgi   = ec.pgi(mat)                    # Product Gini Index
peii  = ec.peii(mat)                   # Product Export Inequality Index
\end{lstlisting}

%% ============================================================
\section{Complexidade Baseada em Patentes (\texttt{patents})}

\begin{lstlisting}[caption={Indicadores de patentes (input: matriz patente x tecnologia)}]
ease_rec = ec.ease_of_recombination(mat_pat)
mod_comp = ec.modular_complexity(mat_pat)
\end{lstlisting}

%% ============================================================
\section{Din\^{a}mica Temporal (\texttt{dynamics})}

\begin{lstlisting}[caption={Crescimento e rastreamento de entrada/sa\'{i}da}]
growth  = ec.growth_rates(df, "regiao", "setor", "empregos", "ano")
entries = ec.entry_tracking(df, "regiao", "setor", "ano")
exits   = ec.exit_tracking(df, "regiao", "setor", "ano")
\end{lstlisting}

%% ============================================================
\section{Complexity Outlook (\texttt{outlook})}

O \textbf{COI} mede o potencial latente de diversifica\c{c}\~{a}o;
o \textbf{COG} mede o ganho esperado de complexidade:

\[
  \text{COI}_c = \sum_p (1 - M_{cp})\,\omega_{cp}\,Q_p
\]

\begin{lstlisting}[caption={COI e COG}]
coi = ec.coi(mat, phi_d, pci)
cog = ec.cog(mat, phi_d, pci)
\end{lstlisting}

%% ============================================================
\section{Pipeline Unificado (\texttt{compute\_complexity})}

\subsection{Assinatura da Fun\c{c}\~{a}o}

\begin{lstlisting}[caption={Assinatura da fun\c{c}\~{a}o pipeline}]
def compute_complexity(
    df: pd.DataFrame,
    cols: dict,           # {"loc": str, "act": str, "val": str, "time": str}
    method: str = "eigenvector",   # 'eigenvector' | 'reflections' | 'fitness'
    threshold: float = 1.0,        # limiar para Mcp
    compute_coi_cog: bool = True,
    iterations: int = 20,          # para reflections/fitness
) -> pd.DataFrame:
    """
    Retorna o DataFrame original enriquecido com as colunas:
    rca, mcp, diversity, ubiquity, eci, pci,
    density, distance, [coi, cog]
    """
\end{lstlisting}

\subsection{Par\^{a}metros}

\begin{center}
\renewcommand{\arraystretch}{1.3}
\begin{tabular}{llp{7cm}}
\toprule
\textbf{Par\^{a}metro} & \textbf{Tipo} & \textbf{Descri\c{c}\~{a}o} \\
\midrule
\texttt{df}              & \texttt{DataFrame} & Tabela no formato longo \\
\texttt{cols}            & \texttt{dict}      & Mapeamento de colunas (ver Se\c{c}\~{a}o~2.3) \\
\texttt{method}          & \texttt{str}       & M\'{e}todo de complexidade \\
\texttt{threshold}       & \texttt{float}     & Limiar do Mcp (padr\~{a}o: 1.0) \\
\texttt{compute\_coi\_cog} & \texttt{bool}    & Calcular COI e COG \\
\texttt{iterations}      & \texttt{int}       & Itera\c{c}\~{o}es para m\'{e}todos iterativos \\
\bottomrule
\end{tabular}
\end{center}

%% ============================================================
\section{Exemplos Completos}

\subsection{Exemplo 1 --- An\'{a}lise Regional de Empregos}

\begin{lstlisting}[caption={An\'{a}lise completa de emprego regional}]
import sys
sys.path.insert(0, "/caminho/para/econcomplex")

import pandas as pd
import numpy as np
import econcomplex as ec

# Gerar dados sinteticos (ou carregar CSV real)
df = ec.make_sample_data(n_locs=50, n_acts=30, seed=42)
# Colunas: loc, act, val

# Pipeline completo
resultado = ec.compute_complexity(
    df,
    cols={"loc": "loc", "act": "act", "val": "val"},
    method="eigenvector",
    compute_coi_cog=True,
)

# Resumo dos indicadores gerados
print(resultado.columns.tolist())
# ['loc','act','val','rca','mcp','diversity','ubiquity',
#  'eci','pci','density','distance','coi','cog']

# Top 10 regioes por ECI
top_eci = (resultado[["loc","eci"]]
           .drop_duplicates()
           .sort_values("eci", ascending=False)
           .head(10))
print(top_eci)

# Salvar
resultado.to_csv("resultado_regional.csv", index=False)
\end{lstlisting}

\subsection{Exemplo 2 --- Exporta\c{c}\~{o}es com Dados COMEX}

\begin{lstlisting}[caption={An\'{a}lise de exportacoes com dados COMEX/MDIC}]
import sys
sys.path.insert(0, "/caminho/para/econcomplex")

import pandas as pd
import econcomplex as ec

# Carregar base COMEX agregada por municipio x NCM x ano
comex = pd.read_parquet("comex_mun_ncm.parquet")
comex_2023 = comex[comex["ano"] == 2023].copy()

# Limpeza
comex_2023 = comex_2023.dropna(subset=["co_mun","co_ncm","vl_fob"])
comex_2023 = comex_2023[comex_2023["vl_fob"] > 0]

# Pipeline
resultado = ec.compute_complexity(
    comex_2023,
    cols={"loc": "co_mun", "act": "co_ncm", "val": "vl_fob"},
    method="eigenvector",
    compute_coi_cog=True,
)

# Salvar
resultado.to_csv("complexidade_exportacoes_2023.csv", index=False)
print("ECI medio:", resultado["eci"].dropna().mean().round(4))
\end{lstlisting}

\subsection{Exemplo 3 --- Painel Temporal RAIS (2010--2022)}

\begin{lstlisting}[caption={Pipeline com dados em painel --- multiplos periodos}]
import sys
sys.path.insert(0, "/caminho/para/econcomplex")

import pandas as pd
import econcomplex as ec

rais = pd.read_parquet("rais_municipio_cnae_2010_2022.parquet")

# Pipeline temporal -- itera sobre cada ano automaticamente
painel = ec.compute_complexity(
    rais,
    cols={
        "loc":  "municipio",
        "act":  "cnae2d",
        "val":  "vinculos",
        "time": "ano",
    },
    method="eigenvector",
    compute_coi_cog=False,
)

# Evolucao do ECI mediano
eci_trend = (painel[["ano","municipio","eci"]]
             .drop_duplicates()
             .groupby("ano")["eci"]
             .median()
             .round(4))
print(eci_trend)

painel.to_csv("painel_complexidade_2010_2022.csv", index=False)
\end{lstlisting}

%% ============================================================
\section{Refer\^{e}ncias Bibliogr\'{a}ficas}

\begin{description}[leftmargin=1.5em, style=nextline]

  \item[Balassa (1965)]
    Balassa, B. (1965). Trade liberalisation and ``revealed'' comparative
    advantage. \textit{The Manchester School}, 33(2), 99--123.

  \item[Hidalgo \& Hausmann (2009)]
    Hidalgo, C.A., \& Hausmann, R. (2009). The building blocks of economic
    complexity. \textit{PNAS}, 106(26), 10570--10575.

  \item[Balland (2017)]
    Balland, P.-A. (2017). \textit{EconGeo: Computing Key Indicators of the
    Geography of Innovation}. R package.

  \item[Vargas Sep\'{u}lveda (2020)]
    Vargas Sep\'{u}lveda, M. (2020). \textit{economiccomplexity: Computational
    Methods for Economic Complexity}. R package.

  \item[The Growth Lab at Harvard (2020)]
    The Growth Lab at Harvard University (2020). \textit{py-ecomplexity:
    Economic Complexity in Python}. GitHub.

  \item[Tacchini et al.\ (2012)]
    Tacchini, E., Loreto, V., Pietronero, L. (2012). A New Metrics for
    Countries' Fitness and Products' Complexity.
    \textit{Scientific Reports}, 2, 723.

\end{description}

%% ============================================================
\section{Licen\c{c}a e Cita\c{c}\~{a}o}

Esta biblioteca \'{e} distribu\'{i}da sob a licen\c{c}a MIT. Para uso em
trabalhos acad\^{e}micos, sugerimos citar os pacotes originais nas vers\~{o}es
correspondentes aos indicadores utilizados.

\begin{lstlisting}[language={}, caption={Verificar vers\~{a}o da biblioteca}]
import econcomplex as ec
print(ec.__version__)   # 0.1.0
\end{lstlisting}

\end{document}
